\documentclass[11pt]{article}

\usepackage[margin=1.1in]{geometry}
\usepackage{amsmath,amssymb,amsthm}
\usepackage{booktabs}
\usepackage{enumitem}
\usepackage{xspace}
\usepackage{xcolor}
\usepackage[colorlinks=true,linkcolor=blue!50!black,citecolor=blue!50!black,urlcolor=blue!50!black]{hyperref}

% ---------------------------------------------------------------------------
% Macros (shared with multi-table-design.tex, circuit-wiring-design.tex)
% ---------------------------------------------------------------------------
\newcommand{\F}{\mathbb{F}}
\newcommand{\Ftwo}{\F_2}
\newcommand{\eq}{\mathsf{eq}}
\newcommand{\nnz}{\mathsf{nnz}}
\newcommand{\ip}[2]{\langle #1, #2 \rangle}
\newcommand{\code}[1]{\texttt{#1}}
\newcommand{\settled}[1]{\medskip\noindent\textbf{SETTLED} (#1).\xspace}

\newtheorem{lemma}{Lemma}
\newtheorem{remark}{Remark}
\theoremstyle{definition}
\newtheorem{definition}{Definition}

\title{Folding the Verifier's Matrix Work:\\ Deferred Claims, Accumulation, and How Recursion Closes\\[0.6ex]
\large Design document --- working draft v0.1}
\author{Flock}
\date{July 31, 2026}

\begin{document}
\maketitle

\begin{abstract}
Flock's lincheck verifier spends most of its time on a single quantity: a
bilinear form in the base matrices $\alpha \hat A_0 + \hat B_0$ at a rank-1
tensor point. It is $O(\nnz)$ --- ${\sim}21$M nonzeros for BLAKE3 --- and
${\sim}84\%$ of verify. This document designs its removal from the verifier:
the verifier stops evaluating it and instead \emph{emits} it as a claim, and
claims about the (registry-static) matrices are \textbf{folded} into a
constant-size accumulator that somebody discharges once, at the end.
Natively this is not a speedup and is not meant to be --- folding $k$ claims
costs the same $k\cdot\nnz$ as checking them. What it buys is an
\textbf{asymmetry}: the fold's prover pays $O(\nnz)$, its verifier pays
$O(\kappa)$ and reads no matrix at all. That is worth nothing outside a
circuit and everything inside one, because arithmetising the matrix work
directly is \emph{$\nnz$-preserving} --- the gadget's matrix \emph{is} the
matrix it evaluates --- so a recursion circuit that pays it inline can never
verify a proof of itself. Deferral plus folding is what makes the fixed
point close. Everything here is implemented on the \code{folding\_verifier}
branch and composed with the circuit layer on \code{recursion\_circuit},
where a native merge node verifies two circuit proofs with their matrix work
deferred and folds both children into one accumulator.
\end{abstract}

\tableofcontents

% ===========================================================================
\section{The problem}\label{sec:problem}

\subsection{What the verifier actually spends its time on}

The union lincheck's final consistency check is, per boolean type $t$,
\[
  \text{target} \;=\; \sum_t P_t\,w_t\bigl(\alpha\ip{W_t}{A_0^{(t)}} +
  \ip{W_t}{B_0^{(t)}}\bigr) \;+\; \sum_t P_t\,\beta_t\,W^{\text{col}}_t[\text{pin}],
\]
where the weight $W_t = W^{\text{row}}_t \otimes W^{\text{col}}_t$ with
\[
  W^{\text{row}}_t = \lambda(z_{\text{skip}}) \otimes \eq(x_{\text{inner\_rest}}),
  \qquad
  W^{\text{col}}_t = z_{\text{partial}} \otimes \eq(rr),
\]
and $w_t, P_t, \beta_t$ are verifier-computable scalars (the slot prefix
weights from the row point and the bound column point, and the const-pin
challenge). Everything the verifier does \emph{except} those bilinear forms
is succinct.

\settled{measured, M30} The split is stark. Of the lincheck verify, the
matrix work is essentially all of it: union-single $9.28/9.36$\,ms,
union-two $19.26/19.34$, mixed (SHA-256 + BLAKE3) $10.15/10.23$. What
remains succinct across the whole proof is ${\sim}1.9$\,ms, and it is
dominated by the \emph{opening}, not the lincheck, whose succinct part is
$0.08$\,ms.

\subsection{Why it cannot simply be arithmetised}\label{sec:nnzpreserving}

The obvious move for recursion --- express the matrix evaluation as gates
--- fails for a structural reason worth stating precisely, because it is the
whole justification for everything that follows.

The comb is a \emph{linear} map, $\text{comb}[c] = \sum_r
\eq_{\text{inner}}[r]\cdot M[r,c]$. In the element ($\F$-native) R1CS a
linear map costs almost no multiplications and almost no area: one
constraint per output, with the coefficients living in the $A_0$ matrix. Its
entire cost appears as \textbf{matrix nonzeros} --- and the gadget computing
it has matrix $M^{\mathsf T}$, so
\[
  \nnz(\text{gadget}) \;=\; \nnz(M).
\]
Arithmetising is $\nnz$-\emph{preserving}. Now impose self-recursion: the
circuit $R$ must evaluate the comb of every table in the proof it verifies,
which is a proof of $R$, so
\[
  \nnz(\text{comb gadget}) \;=\;
  \nnz(\text{Merkle}) + \nnz(\text{FS}) + \nnz(\text{comb gadget}) + \cdots,
\]
which has no solution. Each recursion level adds the \emph{other} tables'
$\nnz$ rather than converging.

\begin{remark}[Constant factors do not help]
Structural factorisation --- exploiting that a composite table is $k$
aligned repeats of a base block, so its comb costs the base block's $\nnz$
rather than $k$ times it --- is a large and worthwhile win natively
(measured $15.5\times$ on the Merkle composite's verify). It does not change
the recurrence above: it shrinks the increment, never making it zero. The
requirement is \emph{sublinearity} in $\nnz$, which only a different
mechanism provides.
\end{remark}

% ===========================================================================
\section{Deferral}\label{sec:defer}

\settled{implemented} The verifier stops discharging the matrix work and
returns it.

\code{lincheck::MatrixAssertion} carries exactly the residue: $\alpha$, the
row weight's $(z_{\text{skip}}, x_{\text{inner\_rest}})$, the column
weight's $(z_{\text{partial}}, rr)$, the pin $\beta$'s, and the target the
weighted bilinear sum must equal. \code{verify\_union\_deferred} returns it
alongside the witness claim; \code{verify\_union} is that followed by
\code{MatrixAssertion::check}, so every existing caller is unchanged and the
byte-identity fixtures still pass.

\paragraph{Why this is a refactor and not a protocol change.} The combs are
needed nowhere before the final consistency check, and building them touches
the challenger not at all. Lifting them out therefore leaves the transcript
byte-identical --- which is what let the split land under the existing
anchors.

\paragraph{The reported evaluations.} The verifier's final check is one
equation in $2T$ unknowns, so it cannot recover the individual bilinear
forms. The prover reports them: \code{LincheckProof::matrix\_evals} carries,
per boolean type, the \textbf{unscaled} pair
$\bigl(\ip{W_t}{A_0}, \ip{W_t}{B_0}\bigr)$. The verifier checks they
reproduce the target using only its own scalars, then hands them out as
claims.

\settled{A and B stay separate} Only their $\alpha$-combination appears in
the target, and $\alpha$ is per-proof --- so a claim about
$\alpha A_0 + B_0$ would name a \emph{different polynomial in every proof}
and could never fold with its neighbours. Claims about $A_0$ and $B_0$
individually name registry-static matrices and fold forever.

\paragraph{Both classes.} The element class defers identically
(\code{element\_r1cs::union::verify\_deferred} $\to$
\code{ElementAssertion}). Element combs are small, which is beside the
point: \S\ref{sec:nnzpreserving} does not care about size. Element weights
are simpler --- plain $\eq \otimes \eq$, since the element PIOP has no
univariate skip and hence no $\lambda$ factor.

\begin{remark}[$z_{\text{eval}} = 0$]
The element relation is $\text{running} = \widehat{\text{Comb}}(r_{\text{col}})
\cdot z_{\text{eval}}$. Stating the assertion by dividing out
$z_{\text{eval}}$ is wrong exactly when it is zero (an all-dummy instance),
where the relation constrains the matrices not at all --- an honest empty
proof gets rejected. The assertion therefore carries both sides and checks
$\widehat{\text{Comb}} \cdot z_{\text{eval}} = \text{running}$: no special
case, no inversion.
\end{remark}

% ===========================================================================
\section{The fold}\label{sec:fold}

A claim is $\sum_{r,c} \text{row}(r)\,\text{col}(c)\,M[r,c] = v$ for
structured weights. \code{matrix\_fold::Weight} is one shape,
$\text{low} \otimes \eq(\text{point})$, covering both what the lincheck
produces ($\lambda \otimes \eq$, $z_{\text{partial}} \otimes \eq$) and what a
folded claim carries (plain $\eq$) --- which is what lets a fold's output
feed straight back in.

\subsection{Two sumchecks, in this order}

$k$ claims about one $M$ fold to one:

\begin{enumerate}[nosep]
  \item \textbf{Column.} With $\text{comb}_i(c) = \sum_r
        \text{row}_i(r) M[r,c]$, each claim reads
        $\sum_c \text{col}_i(c)\,\text{comb}_i(c) = v_i$. A degree-2
        sumcheck on the $\lambda$-combination binds every claim to one
        shared $\rho_{\text{col}}$; the prover then sends the $k$ bridge
        values $\text{comb}_i(\rho_{\text{col}})$.
  \item \textbf{Row.} Because $\rho_{\text{col}}$ is now shared, every
        bridge value reads against the \emph{same}
        $h(r) = \hat M(r, \rho_{\text{col}})$ --- so the row side collapses
        to a single $\mu$-combined sumcheck landing on
        $\hat M(\rho_{\text{row}}, \rho_{\text{col}})$.
\end{enumerate}

The output is a plain evaluation claim, which folds again by the same
machinery.

\settled{the order is load-bearing} The $k$ per-claim marginals are the
dominant $k\cdot\nnz$ cost, and \emph{which kind} they are is decided by
which phase runs first. Column-first makes them \textbf{column} marginals,
which is exactly what \code{LincheckCircuit::fold\_split} computes with each
type's tuned walker/CSC kernel. Row-first would make them row marginals,
which no tuned kernel provides. Only the single shared marginal of step 2
touches the raw matrix.

\begin{remark}[Why fresh $\mu$]
The bridge values are prover-supplied. Reporting them
$\lambda$-consistently but individually wrong shifts the column sumcheck's
target, so the error lands in the output claim --- where the root discharge
catches it. That is the accumulation property: \textbf{if any input claim is
false, the output claim is false with overwhelming probability.}
\end{remark}

\subsection{Nothing is committed}

The claims are about \emph{public, registry-static} polynomials, so the
accumulated claim at the root is discharged by direct evaluation in the
clear. No sparse-matrix commitment, no preprocessing, no verification key
--- which is the SPARK-style delegation this design replaces.

\subsection{Sumcheck prover cost}

\settled{measured, BLAKE3 $k_{\log}=14$, 21M nonzeros, 14 threads} The fold
$+$ discharge went $66.6 \to 10.8$\,ms across three changes, each with a
distinct cause: interleaving the $k$ column weights so a nonzero costs one
contiguous $k$-wide read instead of $k$ scattered ones, and parallelising
both generic passes ($66.6 \to 25.1$); the phase order above, putting the
$k\cdot\nnz$ work on the tuned kernels ($25.1 \to 15.9$); and a one-pass
\code{fold\_split} ($15.9 \to 10.8$).

That last one is worth recording: the default \code{fold\_split} walked the
nonzeros \emph{twice} (\code{fold\_alpha\_batched} at $\alpha = 0$ and
$\alpha = 1$) to separate the two marginals. Both real implementations
already form them apart internally and only mix at the end, so keeping them
apart is strictly \emph{less} work --- the same walk without the per-column
multiply. Measured $2.258 \to 1.050$\,ms, faster than a single
$\alpha$-batched fold.

What remains is near the floor: $k\cdot\nnz$ for the per-claim marginals
(${\sim}4.2$\,ms), one row marginal per fold (${\sim}2.6$\,ms, the direction
no tuned kernel provides), and the discharge (${\sim}2.8$\,ms, which a merge
node does not pay --- the root discharges once for the whole tree).

% ===========================================================================
\section{Accumulation}\label{sec:accum}

\code{aggregate::Accumulator} holds one $(A_0, B_0)$ claim pair per type,
in two groups --- boolean and element --- tied to a registry digest.

\settled{the key is $(\text{registry}, \text{type}, \text{matrix})$} All
three components separate claims that name different polynomials, and
folding across any of them would be meaningless. The digest covers the
first; the class split the second; $A$/$B$ the third. An accumulator from a
different registry is rejected, not silently folded.

\paragraph{How many claims per fold.} Within one accumulator:

\begin{center}
\begin{tabular}{lccl}
\toprule
node & inherited & fresh & fold \\
\midrule
leaf over two base proofs & 0 & 2 & $2\to1$ \\
$2\to1$ merge of two recursive proofs & 2 & 2 & $4\to1$ \\
$1\to1$ self-recursion & 1 & 1 & $2\to1$ \\
\bottomrule
\end{tabular}
\end{center}

\emph{Fresh} is what succinctly verifying a child emits; \emph{inherited} is
the accumulator the child already carried, which the merge cannot discharge
(that is the $O(\nnz)$ it is avoiding) and so carries forward.

\paragraph{Constant size.} After a fold, a claim's weights are plain $\eq$,
so it is $(\rho_{\text{row}}, \rho_{\text{col}}, v)$ --- $2\kappa+1$ field
elements. The accumulator is independent of tree depth and of how many
proofs have been folded in: depth 1 and depth 20 carry the same few
kilobytes. That is what makes unbounded aggregation possible.

\paragraph{An API that cannot be misused.} \code{verify\_aggregate} runs
each assertion's \code{check\_reported} itself, so the step tying a proof's
reported matrix work \emph{to that proof} cannot be forgotten by a caller.

% ===========================================================================
\section{Merge composition, and why recursion closes}\label{sec:merge}

A recursive proof is $\pi = (p, \alpha)$: \emph{$p$ attests to the
statement, assuming the accumulated claims $\alpha$ are true.}

Merging $\pi_1, \pi_2$: the merge circuit succinctly verifies both children
(emitting fresh claims $\beta^1, \beta^2$), folds
$\{\alpha^1, \alpha^2, \beta^1, \beta^2\}$ per (type, matrix), and outputs
the folded $\alpha$ as a public output.

\begin{lemma}[The invariant]
$\pi = (p, \alpha)$ is valid iff $p$ succinctly verifies and $\alpha$ is
true; and the merge preserves it.
\end{lemma}

Verifying child $p_1$ leaves two debts --- the one it just created
($\beta^1$, its own lincheck's matrix evaluation) and the one it arrived
with ($\alpha^1$, everything beneath it). Folding all four covers the entire
subtree. Neither can be dropped: discharging $\alpha$ costs the $O(\nnz)$
the design exists to avoid, and ignoring $\beta$ lets the child's lincheck
go unchecked. The debt telescopes, so whoever discharges the root's $\alpha$
retroactively validates every lincheck in the tree.

\paragraph{Where the fold transcript lives.} It is \textbf{witness} to the
merge circuit. The fold's prover runs natively, before the circuit exists;
its transcript enters as witness, is hashed by the circuit's own FS gadget
to derive $\lambda, \mu, \rho$, and is \emph{consumed} there. It never
appears in the output proof --- the verifier of $\pi$ checks $p$, which
already attests the fold was replayed correctly. That is what preserves the
shape: proof plus accumulator in, proof plus accumulator out.

Folding \emph{outside} and appending the transcript does not close: the
proof grows by a fold transcript per level, unless the next level's circuit
consumes it --- which is this design again.

\paragraph{Why $\alpha$ and $\beta$ fold together.} They are the same kind
of object, differing in weight shape ($\beta$ is $\lambda \otimes \eq$
fresh from a lincheck; $\alpha$ is plain $\eq \otimes \eq$ from a previous
fold), in provenance ($\beta$ is derived now from the child's challenges,
$\alpha$ arrives as public state), and in scope ($\beta$ is the child's own
lincheck, $\alpha$ is everything beneath it). What they share is the matrix
they name --- under self-recursion every proof is of the same circuit, hence
the same registry.

\begin{remark}[Two registries, not one]
At level $\ge 2$ they may name \emph{different} registries: the merge
circuit verifying two \emph{base} proofs emits $\beta$ about the application
circuit's matrices, while at the next level $\beta$ is about the merge
circuit's own and the inherited $\alpha$ is still about the application's.
The groups fold independently and the slot set stabilises at two registries
--- so the accumulator stays constant-size, roughly double a single
registry's.
\end{remark}

\paragraph{Cost.} In-circuit, each fold verifier is $2\kappa$ sumcheck
rounds ($\kappa$ row $+$ $\kappa$ column: the matrix is
$2^\kappa \times 2^\kappa$, square because $C = I$), ${\approx}38$ for the
Merkle composite, plus a few weight evaluations. Its transcript is $4\kappa$
field elements (${\sim}1.2$\,KB), so absorbing it costs ${\sim}20$ BLAKE3
compressions against a ${\sim}10$k budget. The recurrence's increment drops
from ${\sim}42$M to a few hundred constraints, and converges.

\begin{remark}[Two phases, not one sumcheck]
The full matrix domain is $2^{2\kappa}$ --- $2^{38}$ for the Merkle
composite --- and nothing can materialise it. Splitting into row and column
phases means the prover only ever touches length-$2^\kappa$ vectors, so
sparsity is exploited by the \emph{structure of the reduction} rather than
by a sparse-polynomial commitment.
\end{remark}

% ===========================================================================
\section{What this does \emph{not} buy}\label{sec:notbuy}

\settled{measured, and it retracts an earlier claim} Native batch
verification is \textbf{not} faster --- it is several times slower. At
$N = 4$ BLAKE3 proofs: $9.4$\,ms of ordinary verifies against $71$\,ms of
deferred-verify-plus-fold (later $23$\,ms after the fold optimisations, still
slower).

The reason is structural. Folding $k$ claims costs $k \cdot \nnz$, because
the fold builds one marginal per claim --- exactly what checking those $k$
claims directly costs --- with the discharge on top. No random-linear-
combination avoids it: the combined weight is a sum of $k$ rank-1 terms, so
every nonzero still needs $k$ multiplications. \emph{$k$ claims at $k$
distinct points cost $k$ passes over the matrix, full stop.}

What survives is the asymmetry, and it is worth only one thing: the fold's
\textbf{prover} pays the $k\cdot\nnz$; its \textbf{verifier} pays
$O(\kappa)$ and reads no matrix. That moves the matrix work from inside a
circuit --- where it is $\nnz$-preserving and the fixed point cannot close
--- to a native prover, where it is ordinary.

% ===========================================================================
\section{Implementation status}\label{sec:status}

On \code{folding\_verifier} (downstream of \code{multitable}), composed with
the circuit layer on \code{recursion\_circuit}:

\begin{itemize}[nosep]
  \item \code{lincheck::MatrixAssertion} $+$ \code{verify\_union\_deferred};
        \code{verify\_union} is the composition. Byte fixtures re-pinned for
        the \code{matrix\_evals} report.
  \item \code{element\_r1cs::union::verify\_deferred} $+$
        \code{ElementAssertion} --- both classes defer.
  \item \code{matrix\_fold}: \code{Weight}, \code{MatrixClaim},
        \code{FoldMatrix} (spanning $\Ftwo$ supports and $\F$-coefficient
        matrices), \code{prove\_fold} / \code{verify\_fold},
        \code{bilinear} / \code{check\_direct}.
  \item \code{aggregate}: \code{Accumulator} (boolean $+$ element groups,
        registry digest), \code{prove\_aggregate\_classes} /
        \code{verify\_aggregate\_classes}, \code{discharge}.
  \item Deferred verify entries at proof level, for the union and (merged
        only) for circuits.
  \item A native merge node: two circuit proofs verified deferred, both
        children folded into one accumulator, discharged once
        (\code{a\_merge\_node\_folds\_two\_circuit\_proofs}).
\end{itemize}

\paragraph{Tests worth knowing about.} The differential oracles are the
load-bearing ones: an honest fold's \emph{output} is itself true
(\code{check\_direct} against the real matrix), so an accumulator can be
believed at the root; the reported evaluations recombine to the comb the
verifier used to build itself, so deferring changed what is \emph{computed}
and not what is \emph{checked}; and both transports carry the same claim set
while their openings differ.

\begin{remark}[A recurring test-writing mistake]
Twice I asserted that a tampered input makes the \emph{prover's}
accumulator false. It does not: \code{prove\_fold} computes its output from
the real matrix, so the prover's claim is honest regardless of what the
input claims said. It is the \emph{verifier's} accumulator that is
conditional on them. Relatedly, a tampered fold may be caught by outright
rejection \emph{or} by a false accumulator --- both sound, so an assertion
must accept either, and the property is always about the verifier's side.
\end{remark}

% ===========================================================================
\section{Open questions}\label{sec:open}

\begin{enumerate}[nosep]
  \item \textbf{Element-only proofs over the merged transport} ---
        \code{prove\_batched\_padded\_with\_precomputed} asserts at least one
        ring-switched claim, so a proof with no boolean class cannot use it.
        Small and isolated.
  \item \textbf{A row-marginal kernel per table type} would close the last
        generic pass in the fold (${\sim}2.6$\,ms of ${\sim}10.8$). Worth it
        only if merge-prover time becomes the constraint --- a merge also
        generates a full proof, which will dwarf it.
  \item \textbf{The jagged remainder} (see the circuit-wiring document): the
        merged transport's verifier delegates its jagged-composed weight
        evaluation via the Frobenius assist, which is column-linear. Now
        that element and circuit proofs ride the merged path, it is no
        longer avoidable by choosing a transport.
  \item \textbf{Structural factorisation of the base blocks} --- a constant
        factor, not a fix (\S\ref{sec:nnzpreserving}), but it cheapens the
        fold's prover, which every merge node pays.
\end{enumerate}

\end{document}
